The most fundamental limitation of the system is the structural mismatch between a fixed DSL and non-deterministic Python execution.
The plan-and-execute agent is constrained to a finite, pre-specified operator set, while the iterative analytics agent is free to construct arbitrary Python logic at runtime.
Even when both agents arrived at numerically comparable answers, the intermediate computations differed enough that the validator rejected valid results, producing a high rate of false negatives.
This is not a failure of calibration but a category error: a DSL interpreter cannot faithfully reconstruct the behavior of a free-form code executor.

A related limitation is the scalability ceiling of the DSL itself.
Each new dataset quirk --- multi-table joins, provenance tracking across ticker files, edge cases in aggregation --- required the addition of new operators, causing the plan-and-execute system prompt to grow monotonically.
Because the DSL specification must be presented in full for the agent to reason over it, there is no clean way to modularize or lazy-load operator definitions.
In practice, I observed that larger prompts degraded both response latency and reasoning quality as the model's attention was distributed across an increasingly large instruction surface.

Multi-dataset joins were a persistent source of failure in the plan-and-execute agent.
The financial indicators dataset required concatenating per-ticker CSVs while preserving provenance, and early operator implementations produced malformed outputs due to key collisions, silent column drops, and ambiguous duplicate handling.
Although I extended the DSL with join-aware operators to address these cases, the fixes were reactive rather than principled --- each new query pattern risked exposing a gap in the operator set.

Finally, the validator's notion of correctness was underspecified in cases where both agents reached reasonable but numerically distinct answers due to differing assumptions --- for example, whether to drop or impute \texttt{NaN} values, which join key to use as the primary index, or whether to normalize before aggregating.
The system had no mechanism for adjudicating these disagreements, so structurally valid but approach-divergent outputs were indistinguishable from genuine errors.
